% !TeX root = ../机器学习笔记.tex
% 由主文档通过 \input 载入；请编译主文档。
\setcounter{chapter}{13} % 保留原书章号。
\chapter{聚类分析}
\label{chap:ml-clustering}

聚类并非从无标签数据中自动恢复一个唯一答案，而是按照距离、密度或连接关系建立分组。先明确什么结构有用，再选择与该结构相适应的算法。

\section{聚类任务与结构假设}

硬聚类把每个样本分入一个簇，软聚类则分配多个归属权重。簇标签只是名称，交换编号不改变分组，因此有参考标签时也不能直接比较编号相等的比例。一个有效聚类目标通常同时考虑簇内紧密和簇间分离，但“紧密”的定义依赖输入表示与距离。

例如两条弯曲带状点云可能各自连通却不围绕单一中心，按均值距离切分就未必恢复这两条带。不同粒度也可有不同的合理答案。评价前必须明确所关心的是几何压缩、群体解释还是后续任务性能。

\section{均值聚类与交替最小化}

给定 $K$ 个中心 $\mu_1,\ldots,\mu_K$，$K$ 均值目标为 $\sum_i\|x_i-\mu_{c_i}\|^2$。Lloyd 迭代交替执行：固定中心，将样本分配给最近中心；固定非空簇，将中心更新为簇内均值。初始化可随机选样本，多次初始化比较最终目标；并列分配采用固定规则，空簇需要保留旧中心或用明确的重新初始化办法处理。

分配步骤逐样本最小化相应平方距离；更新步骤由均值的最小二乘性质最小化每簇损失，因此标准非空簇更新使目标不增，目标又以零为下界，故目标值收敛。这不单独证明参数唯一收敛或达到全局最优。常用停止条件是分配不再变化、目标改变量很小或达到上限。

\section{层次聚类与连接准则}

凝聚式层次聚类（hierarchical clustering）从每点一个簇开始，反复合并距离最近的两簇。单连接用最小跨簇点距，全连接用最大跨簇点距，平均连接用平均跨簇点距。保留每步合并及其距离便得到树状图，可在某个层次截断为指定簇数。

单连接容易通过一串近邻点把两个主体连起来，形成链状效应；全连接更强调簇的直径，平均连接则折中。没有一个连接准则能对所有形状都代表正确语义。常见朴素实现还需要维护成对距离，规模较大时应考虑存储与计算代价。

\section{密度聚类与谱聚类}

DBSCAN（density-based spatial clustering of applications with noise）给定半径 $\epsilon$ 和最小邻居数 $m$，将半径邻域内含足够样本的点视为核心点，由核心点之间的密度连通关系形成簇；靠近核心点的非核心点可作为边界点，其他点标为噪声。邻居数是否含自身必须约定。边界点同时邻接不同簇时可能有分配歧义，而密度相差很大时单一参数未必合适。

谱聚类（spectral clustering）先构造对称非负相似矩阵 $W$，置 $D_{ii}=\sum_jW_{ij}$、$L=D-W$。恒等式
$v^\top Lv=\tfrac12\sum_{ij}W_{ij}(v_i-v_j)^2$
说明低特征值方向倾向让强相连节点取相近值。选取若干特征向量作为嵌入，再进行聚类。归一化版本需处理零度节点，图构造与嵌入维数都会影响结果。

\section{簇数选择与聚类评价}

对非单样本簇中的样本 $i$，令 $a_i$ 为同簇平均距离，$b_i$ 为与其他各簇平均距离的最小值，轮廓系数为 $s_i=(b_i-a_i)/\max(a_i,b_i)$，退化情况需约定。较大值表示在当前距离下相对更贴近本簇；它可能偏好凸而分离的结构，并非语义正确性的证据。

有参考标签时可用调整兰德指数等不受编号置换影响的指标；无标签时还可看重采样稳定性与下游效用。对多个 $K$ 比较训练平方距离时，目标随 $K$ 增加不会上升，故不能单凭最小训练目标选择簇数。

\section{例题、练习与学习检查}

\begin{example}[一维均值聚类]
数据为 $0,1,4,5$，取 $K=2$，初始中心为 $0,5$。
\end{example}
\begin{solve}
第一次分配得到 $\{0,1\}$ 与 $\{4,5\}$，更新中心为 $0.5,4.5$，目标由初始的二下降为一。再次分配不变，算法停止。这个简单例子可以穷举验证全局最优，但一般数据上的停止不提供同样保证。
\end{solve}

习题：证明均值更新最小化簇内平方距离；构造单连接的链状效应；说明 $K=n$ 时训练目标为何为零而通常缺乏实用分组意义；比较同一数据经特征缩放前后的聚类。





